% Generated draft: requires mathematical and source review.
\documentclass[11pt]{article}
\usepackage[margin=25mm]{geometry}
\usepackage{fontspec}
\setmainfont{Latin Modern Roman}
\usepackage{amsmath,amssymb,mathtools,unicode-math}
\setmathfont{Latin Modern Math}
\usepackage{enumitem,needspace,xurl,hyperref}
\setlength{\parindent}{0pt}
\setlength{\parskip}{5pt}
\setlength{\emergencystretch}{3em}
\newcommand{\N}{\mathbb{N}}
\newcommand{\Z}{\mathbb{Z}}
\newcommand{\Q}{\mathbb{Q}}
\newcommand{\R}{\mathbb{R}}
\newcommand{\C}{\mathbb{C}}
\newcommand{\F}{\mathbb{F}}
\newcommand{\A}{\mathbb{A}}
\newcommand{\PP}{\mathbb{P}}
\newcommand{\E}{\mathbb{E}}
\newcommand{\eps}{\varepsilon}
\newcommand{\dd}{\,\mathrm{d}}
\begin{document}
\section*{UnsolvedMath ID: 800766}
\textbf{Source problem code:} catalogue entry
\subsection{Definitions and mathematical statement}
For a smooth complex projective variety, rational connectedness means that two general points can be joined by a rational curve; uniruledness means that a general point lies on a rational curve. An MRC (maximal rationally connected) fibration is the rational quotient by the equivalence relation generated by chains of rational curves, considered up to birational change of its base. It is almost holomorphic: it is defined near its general fibres, and those fibres are complete. These definitions describe the usual smooth complex projective setting; the supplied excerpt does not specify the ground field or the full class of varieties under consideration.

Let \(X\) be a projective variety in a setting where an MRC fibration
\[
q\colon X\dashrightarrow Y
\]
is defined. Say that this fibration admits a holomorphic representative, with the source \(X\) fixed, if there are an admissible base variety \(Y'\), a birational map \(\beta\colon Y\dashrightarrow Y'\), and an everywhere defined holomorphic map \(f\colon X\to Y'\) such that
\[
f=\beta\circ q
\]
as rational maps. The source does not state the precise admissibility requirements on \(Y'\), such as normality or projectivity.

\textbf{Problem (Matsumura, 2022, Problem 1).} Determine when an MRC fibration can be chosen to have a holomorphic representative by changing the birational model of its base. In particular, determine the conditions on \(X\) under which such \(Y'\), \(\beta\), and \(f\) exist. This is a request for a criterion, rather than an explicitly stated universal existence conjecture; the intended hypotheses on \(X\) remain unspecified in the supplied excerpt.

The supplied motivation reports a projective variety \(Z\), constructed by Fujita, with Picard number \(\rho(Z)=1\), which is uniruled but not rationally connected. Here \(\rho(Z)\) is the dimension of the real vector space of Cartier divisors modulo numerical equivalence. The source states that \(Z\) has no non-trivial fibrations, hence no holomorphic MRC fibration, and that \(Z\) is neither log canonical nor \(\Q\)-factorial. It refers to \([\mathrm{EIM20},\text{ Example }6.5]\) for details. The source author reports not knowing an analogous example among smooth projective varieties or varieties with mild singularities. The meaning of ``mild singularities'' and the precise convention for ``non-trivial fibration'' are not specified. These are supplied contextual claims, not independently verified assertions or a current literature assessment.
\subsection{Short English statement}
Matsumura asks when changing the birational model of the base of an MRC fibration makes the map everywhere holomorphic while keeping the original projective variety fixed. A supplied singular example obstructs general existence, but the source reports no known analogous smooth example in its 2022 account. The ground field, admissible base varieties, and intended singularity hypotheses require clarification.
\subsection{Sources}
\begin{itemize}
\item Local UnsolvedMath registry: permanent ID \(800766\), exact source problem code \texttt{catalogue entry}, title ``Positively Curved Projective Varieties --- Problem 1.'' The supplied record attributes the problem to Matsumura, with proposed year 2022, and labels it open. This label is supplied registry metadata, not an independently verified mathematical status.
\item Matsumura, \emph{Open Problems on Structure of Positively Curved Projective Varieties} (2022), Problem 1. Supplied source URL: \url{https://arxiv.org/abs/2211.09153}. The provenance records extraction as \texttt{arxiv\_latex} and access on 2026-10-16. This access date is supplied provenance; no independent source visit or verification is claimed.
\item The supplied excerpt attributes the singular example to Fujita and refers to \([\mathrm{EIM20},\text{ Example }6.5]\). No full bibliographic entry or URL for \([\mathrm{EIM20}]\), and no separate Fujita reference, is supplied.
\end{itemize}
\end{document}
